\documentclass[11pt]{article}

\usepackage[margin=1in]{geometry}
\usepackage[T1]{fontenc}
\usepackage{lmodern}
\usepackage{microtype}
\usepackage{amsmath,amssymb,amsthm,mathtools}
\usepackage{booktabs}
\usepackage{enumitem}
\usepackage{xcolor}
\usepackage{xurl}
\usepackage[hidelinks]{hyperref}
\usepackage[nameinlink,capitalize,noabbrev]{cleveref}

\newtheorem{theorem}{Theorem}[section]
\newtheorem{proposition}[theorem]{Proposition}
\newtheorem{lemma}[theorem]{Lemma}
\newtheorem{corollary}[theorem]{Corollary}
\newtheorem{definition}[theorem]{Definition}
\newtheorem{remark}[theorem]{Remark}

\newcommand{\Gr}{\operatorname{Gr}}
\newcommand{\A}{\mathcal A}
\newcommand{\U}{\mathcal U}
\newcommand{\V}{\mathcal V}
\newcommand{\Lcal}{\mathcal L}
\newcommand{\Kcal}{\mathcal K}
\newcommand{\Qcal}{\mathcal Q}
\newcommand{\Hcal}{\mathcal H}
\newcommand{\R}{\mathbb R}
\newcommand{\1}{\mathbf 1}
\newcommand{\ind}{\mathbf 1}
\newcommand{\Span}{\operatorname{span}}
\newcommand{\rank}{\operatorname{rank}}
\newcommand{\supp}{\operatorname{supp}}
\newcommand{\cl}{\operatorname{cl}}
\newcommand{\N}{\operatorname N}
\newcommand{\cert}[1]{\hyperref[cert:#1]{\texttt{#1}}}

\title{The Magic Number Conjecture for the Rank-Two $m=4$ Amplituhedron\\
\large A simplicial reduction of global tile relations to finite local data}
\author{Jack Tavakoli\\University of Houston}
\date{September 28, 2026 --- publication supplement candidate}

\begin{document}
\maketitle

\begin{abstract}
For positive external data $Z\in\R^{n\times 6}$, the rank-two $m=4$
amplituhedron $\A_{n,2,4}(Z)$ is the image of $\Gr_{\ge 0}(2,n)$ under
$C\mapsto CZ$.  We prove that every complete multiplicity-one all-$Z$
positroid tiling of $\A_{n,2,4}$ has
\[
M(2,n-6,2)
=\frac{1}{n-3}\binom{n-3}{2}\binom{n-3}{3}
=\frac{(n-3)(n-4)^2(n-5)}{12}
\]
tiles, for every admissible $n\ge 6$.

The proof separates local and global issues.  Let $V_n$ be the real vector
space generated by all eight-dimensional positroid source cells whose
amplituhedron maps are injective for every positive $Z$.  Let $L_n$ be the
span of the geometric codimension-one parent relations, and let $K_n$ be
the kernel of the family of all-$Z$ tile-indicator maps, modulo
semialgebraic sets of dimension at most seven.  A finite rank-two
classification proves that every local parent relation is balanced.  The
main global theorem is
\[
K_n=L_n\qquad(n\ge 6).
\]
Its proof uses signed adjacent-column amalgamation as face maps and zero
insertion as degeneracy maps.  These operations make $V$, $L$, $K$, and
$K/L$ into simplicial real vector spaces.  Normalization then reduces the
infinite problem to finitely many support cores.  The normalized quotient
has dimensions $1,9,19,12,0,0$ in source label numbers $6,\dots,11$ and
vanishes thereafter by a support bound.  Since $K_9=L_9$ is certified
exactly, exactness and faithfulness of simplicial normalization force
$K/L=0$ in every degree.

The proof is computer-assisted only in finite classification and exact-rank
certificates; all matrix computations are over $\mathbb Q$ or $\mathbb Z$
and are accompanied by exact rank certificates, including nonzero minors
and upper-rank certificates where needed.  No
floating-point rank calculation, exhaustive tiling enumeration, or
fixed-$Z$ extrapolation is used.
\end{abstract}

\paragraph{Status of this draft.}
This manuscript presents a computer-assisted proof with an accompanying
exact certificate supplement.  It has not yet undergone external peer
review.  Each computational claim below carries a stable certificate ID;
\Cref{app:certificates} describes the literal data and verification
conditions.  The supplement's \path{CERTIFICATE_INDEX.md} maps these IDs
to files and replay commands.  The supplement should accompany every
circulated version.

\section{Introduction}

The amplituhedron $\A_{n,k,m}$ was introduced as a positive-geometric model
for scattering amplitudes and has since developed a substantial
combinatorial theory.  Karp--Williams--Zhang conjectured the existence,
for even $m$, of a positroid tiling of $\A_{n,k,m}$ of cardinality
\[
M\!\left(k,n-k-m,\frac m2\right),
\]
the MacMahon number counting plane partitions in a
$k\times(n-k-m)\times(m/2)$ box \cite[Conjecture 8.1]{KWZ}.
The assertion that every tiling has this cardinality is formulated in
\cite[Conjecture 3.9]{GalashinLam}.  Parisi,
Sherman-Bennett, Tessler, and Williams proved the conjecture for $m=2$
using T-duality and Parke--Taylor weights \cite{M2Magic}.  For $m=4$,
Even-Zohar, Lakrec, Parisi, Tessler, Sherman-Bennett, and Williams proved
that all iterated BCFW constructions give genuine tilings
\cite{BCFWCluster}; a companion paper exhibited the non-BCFW spurion tile
\cite{ClusterResults}.  Thus the cardinality problem for arbitrary
positroid tilings is strictly stronger than a statement about BCFW
tilings alone.

The purpose of this paper is to settle the first higher-$m$ rank slice:
$k=2$, $m=4$, and arbitrary $n$.  The main result is the following.

\begin{theorem}[Magic Number for $\A_{n,2,4}$]\label{thm:main}
For every $n\ge 6$, every complete multiplicity-one all-$Z$ positroid
tiling $\mathcal T$ of $\A_{n,2,4}$ has
\[
|\mathcal T|
=M(2,n-6,2)
=\frac{1}{n-3}\binom{n-3}{2}\binom{n-3}{3}
=\frac{(n-3)(n-4)^2(n-5)}{12}.
\]
\end{theorem}

The proof is organized around relations among tile indicators.  For each
$n$ we define a free tile space $V_n$, the span $L_n$ of local relations
coming from nine-dimensional positroid parents, and the full all-$Z$
indicator kernel $K_n$.  The local theorem proves that every vector in
$L_n$ has coefficient sum zero.  The global problem is therefore reduced
to proving
\[
K_n=L_n.
\]

At rank two, this equality admits an unexpectedly finite proof.  Adjacent
column amalgamation and zero insertion turn the system into a simplicial
vector space.  Passing to the normalized quotient discards all generators
obtained by inserting loops.  A classification of all rank-two $m=4$
tiles shows that a tile has at most ten nonloop labels.  Consequently no
new normalized generators can occur after a finite degree.  Two exact
finite-rank computations at label numbers $10$ and $11$, together with the
previously certified equality $K_9=L_9$, close the global argument.

\subsection*{Outline}
\Cref{sec:background} fixes conventions.  \Cref{sec:membership} gives the
finite membership-core theorem.  \Cref{sec:local} states the complete
local parent theorem and explains why all local relations are balanced.
\Cref{sec:amalgamation} proves geometric descent under adjacent
amalgamation.  \Cref{sec:simplicial} establishes the simplicial identities.
\Cref{sec:normalization} gives the finite normalized-rank calculation, and
\Cref{sec:global} proves $K_n=L_n$ for all $n$.  The Magic Number formula is
deduced in \Cref{sec:magic}.  Appendices summarize the exact finite
certificates and reproducibility requirements.

\section{Positive Grassmannians, tiles, and indicator relations}
\label{sec:background}

\subsection{Rank-two positroid cells}
A point of $\Gr(2,n)$ is represented by a full-rank $2\times n$ matrix
$C=(c_1,\dots,c_n)$ modulo left multiplication by $\mathrm{GL}_2$.  Write
$\Delta_{ij}(C)$ for the ordered $2\times2$ Pl\"ucker coordinate.  The
totally nonnegative Grassmannian is
\[
\Gr_{\ge0}(2,n)=\{C:\Delta_{ij}(C)\ge0\text{ for }i<j\}.
\]
Its positroid stratification is recalled from Postnikov \cite{Postnikov}.
In rank two a cell is determined by its loops together with cyclically
ordered parallel classes among the nonloop columns.  If $q$ is the number
of nonloop labels and $s$ the number of cyclic parallel classes, then
\begin{equation}\label{eq:rank2dim}
\dim S=q+s-4.
\end{equation}

\subsection{The amplituhedron map}
Fix a positive matrix $Z\in\R^{n\times 6}$, meaning that all increasing
$6\times6$ minors are positive.  The rank-two $m=4$ amplituhedron map is
\[
f_Z:\Gr_{\ge0}(2,n)\longrightarrow \Gr(2,6),\qquad C\longmapsto CZ,
\]
with $CZ$ interpreted by row span.  Its image is $\A_{n,2,4}(Z)$, which has
dimension $8$.

\begin{definition}[All-$Z$ tile]
An eight-dimensional positroid cell $S\subset\Gr_{\ge0}(2,n)$ is an
\emph{all-$Z$ tile source} if $f_Z|_S$ is injective for every positive
$n\times6$ matrix $Z$.  We write
\[
T_S(Z)=f_Z(\overline S)\subset\A_{n,2,4}(Z)
\]
for its closed tile image.
\end{definition}

The all-$Z$ convention is essential throughout.  A cell that is injective
for one chosen positive matrix but not uniformly in $Z$ is not included in
the tile universe used here.

\begin{definition}[Complete multiplicity-one all-$Z$ tiling]
A finite collection $\mathcal T$ of distinct all-$Z$ tile sources is a
\emph{complete multiplicity-one all-$Z$ positroid tiling} if the collection
is independent of $Z$ and, for every positive $n\times6$ matrix $Z$,
\[
\A_{n,2,4}(Z)=\bigcup_{S\in\mathcal T}T_S(Z),
\qquad
f_Z(S)\cap f_Z(S')=\varnothing\quad(S\ne S').
\]
Thus the closed tiles cover the entire amplituhedron, and their open
images are pairwise disjoint.  In particular,
$\sum_{S\in\mathcal T}\ind_{T_S(Z)}=\ind_{\A_{n,2,4}(Z)}$
modulo semialgebraic sets of dimension at most seven; common tile
boundaries are allowed, and no face-to-face condition is imposed.
\end{definition}

\subsection{Relations modulo lower-dimensional sets}
Let $\U_n$ denote the finite set of all all-$Z$ tile sources and put
\[
V_n=\R^{\U_n}=\bigoplus_{S\in\U_n}\R e_S.
\]
For a fixed positive $Z$, two indicator functions are identified when they
differ only on a semialgebraic subset of dimension at most $7$.  Define the
uniform all-$Z$ indicator kernel
\begin{equation}\label{eq:Kn}
K_n=\left\{c\in V_n:
\sum_{S\in\U_n}c_S\ind_{T_S(Z)}=0
\pmod{\dim\le7}
\text{ for every positive }Z\right\}.
\end{equation}
The coefficients $c_S$ are independent of $Y$ and $Z$.

A complete multiplicity-one all-$Z$ tiling $\mathcal T$ defines a binary
vector
\[
\chi_{\mathcal T}=\sum_{S\in\mathcal T}e_S\in V_n.
\]
If $\mathcal T$ and $\mathcal T'$ tile the same amplituhedron, then
\begin{equation}\label{eq:tilingdifference}
\chi_{\mathcal T}-\chi_{\mathcal T'}\in K_n.
\end{equation}
Thus controlling $K_n$ controls all tiling differences without requiring
flip connectivity or enumeration of all tilings.

\section{The complete rank-two membership classification}
\label{sec:membership}

The global argument will use the fact that all tile sources have bounded
nonloop support.  We therefore record the complete rank-two classification.

\begin{theorem}[Membership-core theorem]\label{thm:membership}
Let $S\subset\Gr_{\ge0}(2,n)$ be an eight-dimensional positroid cell.  Let
$q$ be its number of nonloop labels and let
$b_1,\dots,b_s$ be the cyclic parallel-class sizes.  Then $S$ is an
all-$Z$ tile source if and only if either
\begin{equation}\label{eq:membershipcriterion}
\max_i b_i=q-5,
\end{equation}
or
\[
(q;s;b_1,b_2,b_3)=(9;3;3,3,3).
\]
Up to signed dihedral symmetry and insertion of loops, there are exactly
eleven admissible loopless core types: ten satisfying
\eqref{eq:membershipcriterion} and the $(3,3,3)$ spurion.
\end{theorem}

\begin{proof}
By \eqref{eq:rank2dim}, dimension eight gives $q+s=12$, hence
\[
(q,s)=(6,6),(7,5),(8,4),(9,3),(10,2).
\]
If a parallel class $D$ has at most four complementary nonloop labels,
then a nonzero source row functional annihilating the direction of $D$
produces a nonzero image vector supported on at most four fixed external
rows.  Therefore the image plane meets a fixed external $t$-space with
$t\le4$.  The corresponding incidence locus in $\Gr(2,6)$ has dimension
at most $t+3\le7$, so the eight-dimensional cell cannot map injectively to
an open eight-dimensional image.

If no class has exactly five complementary labels and no class has at most
four complementary labels, the five possible pairs $(q,s)$ force only two
remaining profiles: $(2,2,2,2)$ and $(3,3,3)$.  The former is the
four-mass cell and is not all-$Z$ injective; see Galashin--Lam
\cite{GalashinLam} and the exact two-preimage certificate in the
supplementary computation.  The latter is the spurion, whose all-$Z$
injectivity is the published input stated in
\cite[Section 5.1]{ClusterResults}.  The explicit four-mass data are
\cert{MEM-FOUR-MASS}; the spurion inverse fixtures are
\cert{MEM-SPURION}.  The latter fixtures check the inverse implementation;
they do not replace the all-$Z$ injectivity input.

It remains to prove injectivity for the complement-five family.  Delete
loops and cyclically cut so that a distinguished class $D$ is last and
its five complementary labels occur first.  If
$A_1,\dots,A_\ell$ are the successive complementary classes and
$r=|D|$, then $\ell+r=6$ and every open source point admits the chart
\[
C=\begin{pmatrix}
 a_1&\cdots&a_5&0&\cdots&0\\
 t_{j(1)}a_1&\cdots&t_{j(5)}a_5&d_1&\cdots&d_r
\end{pmatrix},
\qquad a_i,d_h>0,\quad t_1<\cdots<t_\ell.
\]
Let $H=\Span(Z_1,\dots,Z_5)$ and
$h(X)=\det(Z_1,\dots,Z_5,X)$.  The first image row lies in $H$, while the
second has positive $h$-value, so $Y\cap H$ is a unique line.  Its positive
ray uniquely determines $a_1,\dots,a_5$ up to common scale.  Put
\[
V_j=\sum_{i\in A_j}a_iZ_i.
\]
The six rows
\[
(V_1,\dots,V_\ell,Z_{D_1},\dots,Z_{D_r})
\]
are independent because their determinant is a positive sum of increasing
$6\times6$ minors of $Z$.  Expanding any second image row of positive
$h$-value in this basis recovers all slopes $t_j$ and all $d_h$ uniquely.
Thus the source row space is uniquely recovered.  This proves all-$Z$
injectivity for all ten complement-five cores.

The finite composition list is then obtained from the five possibilities
$q+s=12$ up to cyclic rotation and reversal.  The eleven admissible core
types are displayed in \Cref{tab:membershipcores}.  Loop insertion and
signed dihedral transport preserve the entire source/image problem, so the
classification holds for arbitrary $n$.
The complete labeled catalogue, its dihedral orbits, and the finite
composition enumeration are supplied as \cert{MEM-CORES}.
\end{proof}

\begin{table}[ht]
\centering
\begin{tabular}{ccl}
\toprule
$q$ & number of loopless cells & admissible class-size cores\\
\midrule
6 & 1  & $(1,1,1,1,1,1)$\\
7 & 14 & $(1,1,1,2,2)$, $(1,1,2,1,2)$\\
8 & 36 & $(1,1,3,3)$, $(1,2,2,3)$, $(1,2,3,2)$, $(1,3,1,3)$\\
9 & 30 & $(1,4,4)$, $(2,3,4)$, $(3,3,3)$\\
10& 5  & $(5,5)$\\
\bottomrule
\end{tabular}
\caption{The eleven admissible loopless membership cores.  The
$(3,3,3)$ profile is the spurion; the other ten are complement-five.
The labeled counts and orbit data are certificate \cert{MEM-CORES}.}
\label{tab:membershipcores}
\end{table}

\begin{corollary}[Support bound and tile-source count]\label[corollary]{cor:supportbound}
Every all-$Z$ tile source has at most ten nonloop labels.  Moreover
\begin{equation}\label{eq:Uncount}
|\U_n|
=\binom n6+14\binom n7+36\binom n8+30\binom n9+5\binom n{10}.
\end{equation}
\end{corollary}

The counts in \eqref{eq:Uncount} are replayed by \cert{MEM-CORES};
\cert{SUPPORT-CUTOFF} checks their finite support range and the resulting
normalized cutoff used below.

\section{The local parent relation theorem}
\label{sec:local}

Let $P$ be a nine-dimensional positroid cell and let
\[
\mathcal F(P)=\{S:\dim S=8,\ S\lessdot P,\ S\in\U_n\}
\]
be its admissible codimension-one source facets.  Define
\[
\mathcal R_{\mathrm{all}}(P)=
\left\{c\in\R^{\mathcal F(P)}:
\sum_{S\in\mathcal F(P)}c_S\ind_{T_S(Z)}=0
\pmod{\dim\le7}
\text{ for every positive }Z\right\}.
\]
Let $L_n\subset V_n$ be the span of all such geometric one-parent
relations, together with their loop and signed-dihedral transports.

\begin{theorem}[Complete one-parent theorem]\label{thm:local}
For every nine-dimensional rank-two positroid parent $P$ and every
$n\ge7$, the space $\mathcal R_{\mathrm{all}}(P)$ is either zero or the
one-dimensional span of a primitive balanced vector.  Up to loop deletion
and dihedral relabeling, the $29$ loopless parent cores split as follows:
\begin{center}
\begin{tabular}{lccc}
\toprule
sector & number of cores & admissible facets & relation space\\
\midrule
collapsed parent image & 12 & 0 & $0$\\
complement-five circuit & 10 & 6 & $\R\varepsilon$, $3\leftrightarrow3$\\
reflection/spurion & 5 & 10 & $\R\varepsilon$, $5\leftrightarrow5$\\
hostile four-mass & 2 & 8 & $0$\\
\bottomrule
\end{tabular}
\end{center}
In particular, every $r\in L_n$ satisfies
\begin{equation}\label{eq:augmentationzero}
\sum_{S\in\U_n}r_S=0.
\end{equation}
\end{theorem}

\begin{proof}[Proof structure]
For rank two, a nine-dimensional parent has $q+s=13$.  Up to cyclic
rotation and reversal this yields $29$ loopless core types.  The source
facets of every core are enumerated exactly and are classified using
\Cref{thm:membership}: each facet is either an all-$Z$ tile, collapsed by
the class-complement obstruction, or the noninjective four-mass type.

For ten cores a distinguished class has exactly five complementary labels.
The positive fiber is an affine interval whose six actual wall functions
have alternating derivatives, obtained from the Gale vector of a positive
$7\times6$ matrix.  This gives a multiplicity-one $3\leftrightarrow3$
relation.  Four further cores admit an orientation-reversing source
reflection paired with an orientation-preserving target action; the ten
facets therefore pair with opposite image signs, giving
$5\leftrightarrow5$.  The remaining balanced core is the repaired spurion
pattern, whose wall-sign sum is zero by an ordered-minor/Pl\"ucker
calculation.

It remains to show that no additional unbalanced relation exists on the
same facet set.  For the six-wall cores, opposite-side open overlaps form
$K_{3,3}$, forcing every coefficient relation to be a scalar multiple of
the known balanced vector.  For each ten-wall core, nine exact rational
open-overlap witnesses form a spanning tree on the ten facets; the nine
edge equations have rank nine and therefore one-dimensional kernel.
For the two hostile cores, eight exact target neighborhoods produce an
$8\times8$ tile-indicator matrix of determinant $1$, so their relation
spaces are zero.  All determinants and membership/nonmembership decisions
are checked by exact replay.  The complete catalogue is
\cert{LOCAL-PARENTS}; its literal primitive vectors and their zero
coefficient sums are \cert{LOCAL-RELATIONS}.  The $45$ overlap witnesses
and five rank-nine matrices are \cert{LOCAL-TEN-WALL}, and the two hostile
matrices are \cert{LOCAL-HOSTILE-8} and \cert{LOCAL-HOSTILE-9}.
Loop insertion and signed dihedral
transport extend the finite core statement to all $n$.
\end{proof}

\begin{remark}[Computer-assisted scope]
The finite core enumeration, the ten-wall spanning-tree certificates, and
the two hostile determinant certificates are exact finite computations.
No floating-point geometry is used.  The symbolic complement-five family
and the transport arguments are uniform in positive $Z$.
\end{remark}

\section{Adjacent amalgamation and geometric descent}
\label{sec:amalgamation}

We now introduce the global structure that makes the arbitrary-$n$ problem
finite.  For a source matrix $C=(c_1,\dots,c_n)$ define ordinary adjacent
amalgamation by replacing two adjacent columns by their sum.  The cyclic
wrap requires the rank-two sign twist.

For $0\le i<n-1$, set
\begin{equation}\label{eq:facesordinary}
d_i(c_1,\dots,c_n)
=(c_1,\dots,c_{i+1}+c_{i+2},\dots,c_n),
\end{equation}
and set
\begin{equation}\label{eq:facewrap}
d_{n-1}(c_1,\dots,c_n)=(c_1-c_n,c_2,\dots,c_{n-1}).
\end{equation}
On the free tile space, retain the resulting source label only when it is
again an admissible eight-dimensional tile source; otherwise define the
image to be zero.

\begin{theorem}[Geometric amalgamation]\label{thm:amalgamation}
For every $n\ge7$, the maps $d_i:V_n\to V_{n-1}$ have the following
properties.
\begin{enumerate}[label=(\roman*)]
\item Each source face has a single exact positroid-cell image.  If source
dimension is preserved and the output is admissible, the map between the
open source cells is a positive rational isomorphism.
\item If an eight-dimensional output is discarded, then it is collapsed,
not four-mass; its target image has dimension at most seven for every
positive downstairs external matrix.
\item The maps descend to all-$Z$ indicator kernels:
\[
d_i(K_n)\subseteq K_{n-1}.
\]
\item They descend to the local relation spans:
\[
d_i(L_n)\subseteq L_{n-1}.
\]
\end{enumerate}
\end{theorem}

\begin{proof}
We give the ordinary last-pair case; all other ordinary positions follow
by signed cyclic transport, and the wrap is the same geometric operation in
twisted cyclic coordinates.

Write the last two source columns as $a,b$ and their amalgamation as
$w=a+b$.  Every output ordered minor is either an unchanged source minor
or
\[
\Delta_{i,w}=\Delta_{i,n-1}+\Delta_{i,n}\qquad(i<n-1).
\]
Since all summands are nonnegative, there is no cancellation; the full
source zero pattern therefore determines a unique output positroid label.
Rank loss on the nonnegative Grassmannian can occur only at the coordinate
two-plane supported on the merged pair.  Indeed, if every minor involving
$w$ vanishes, the nonnegative summands vanish separately, and any surviving
outside column would force all columns to be parallel, contradicting rank
two.

A direct case analysis on loops, same-class amalgamation, and distinct
parallel classes shows that every dimension-preserving nonloop face admits
a rational inverse.  Geometrically, the output direction $w$ lies strictly
between the two surviving old directions; its positive decomposition in
those two rays uniquely recovers the original pair up to source row-space
gauge.  Faces involving a disappearing singleton or same-class pair lower
dimension and are treated only as lower-dimensional outputs.

By \Cref{thm:membership}, a dimension-preserving nonloop face creates a
singleton class, so an inadmissible eight-dimensional output cannot be the
four-mass profile $(2,2,2,2)$.  Every other excluded eight-cell has a class
with at most four complementary nonloops and is therefore collapsed by the
incidence-dimension argument used in \Cref{thm:membership}.

Fix an arbitrary positive downstairs external matrix $W$.  It can be
realized as a positive degeneration of upstairs external matrices
$Z(\epsilon)$ in which the relevant adjacent rows coalesce as
$\epsilon\to0^+$.  One may choose the perturbation so that every increasing
six-minor remains positive for sufficiently small $\epsilon$.

The only source rank-loss limit is the coordinate plane supported on the
merged pair.  Every limiting target from this locus contains the duplicated
external row and therefore lies in a Schubert locus of dimension four.
All collapsed or lower-dimensional output cells contribute target sets of
dimension at most seven.  Away from the finite union of these exceptional
sets and source-boundary/critical images, surviving tile indicators are
locally constant under the degeneration: membership follows from the
rational source isomorphism and the parameter-dependent inverse function
theorem, while closed nonmembership follows by compactness of source
closures.

Now let $r\in K_n$.  On a generic downstairs open neighborhood, every
upstairs tile indicator appearing in $r$ stabilizes to the corresponding
downstairs indicator, with multiplicity retained.  If $d_i(r)$ were
nonzero on that neighborhood, the same nonzero value would persist for all
small positive $\epsilon$, contradicting the all-$Z$ relation defining
$r$.  Hence $d_i(K_n)\subseteq K_{n-1}$.

Finally take a primitive one-parent relation.  The exact image of its
nine-dimensional parent is either nine-dimensional, eight-dimensional, or
smaller.  In the first case the surviving facet relation belongs to the
already classified one-parent kernel of the downstairs parent.  In the
second case all surviving facets have the same eight-dimensional output
label, and the resulting scalar relation must be zero on its nonempty open
image.  In the last case every surviving output is lower-dimensional.
Thus $d_i(L_n)\subseteq L_{n-1}$ without using the global equality to be
proved.
\end{proof}

\section{The twisted simplicial structure}
\label{sec:simplicial}

For $0\le i<n$, let $s_i:V_n\to V_{n+1}$ insert a zero source column after
$c_{i+1}$.  This is loop insertion.  On source closures,
\begin{equation}\label{eq:zeroinsertion}
(j_I C)Z=CZ_I,
\end{equation}
so zero insertion preserves all-$Z$ tile relations and local parent
relations exactly.

\begin{theorem}[Simplicial structure]\label{thm:simplicial}
Set $V_n=L_n=K_n=0$ for $n\le5$.  With simplicial degree $p=n-1$, the
spaces
\[
V_{p+1},\qquad L_{p+1},\qquad K_{p+1},\qquad
\Qcal_p:=V_{p+1}/L_{p+1},\qquad
\Hcal_p:=K_{p+1}/L_{p+1}
\]
form simplicial real vector spaces, and $\Hcal$ is a simplicial subobject of
$\Qcal$.
\end{theorem}

\begin{proof}
At the raw column level, the face and degeneracy maps satisfy
\begin{align}
d_i d_j&=d_{j-1}d_i &&(i<j),\label{eq:simp1}\\
d_i s_j&=
\begin{cases}
s_{j-1}d_i,&i<j,\\
\mathrm{id},&i=j\text{ or }i=j+1,\\
s_jd_{i-1},&i>j+1,
\end{cases}\label{eq:simp2}\\
s_i s_j&=s_{j+1}s_i &&(i\le j).\label{eq:simp3}
\end{align}
For non-wrap faces these are associativity and insertion identities.  The
two endpoint computations involving the wrap are
\[
(x_0-x_{n-1})+x_1=(x_0+x_1)-x_{n-1},
\]
and
\[
(x_0-x_{n-1})-x_{n-2}=x_0-(x_{n-2}+x_{n-1}).
\]
The face--degeneracy identities at the cyclic boundary are checked in the
same signed convention.

It remains to justify truncation: a raw intermediate that has been sent to
zero must never become an admissible eight-cell later on one side of an
identity.  Rank loss and source-dimension drop cannot be repaired by later
faces or zero insertion.  If the discarded intermediate is an
inadmissible eight-cell, \Cref{thm:membership} and the singleton created by
a dimension-preserving nonloop amalgamation imply that it is collapsed,
not four-mass.  Its class-complement obstruction persists under every
subsequent dimension-preserving face: amalgamating away from the obstructing
class decreases the complement, while amalgamating the obstructing class
decreases both its size and total support.  Hence a discarded intermediate
remains discarded.  The raw equalities therefore descend to the truncated
free tile spaces.

By \Cref{thm:amalgamation}, faces preserve $K$ and $L$.  By
\eqref{eq:zeroinsertion}, degeneracies preserve $K$ and $L$.  The simplicial
identities consequently restrict to these subspaces and descend to their
quotients.
\end{proof}

\section{Normalization and the finite quotient}
\label{sec:normalization}

For a simplicial real vector space $X$, let $N(X)$ denote its normalized
chain complex.  We use the standard facts that $N(X)_p$ is canonically
isomorphic to the quotient of $X_p$ by the span of degeneracies from degree
$p-1$, that normalization is exact, and that it detects the zero simplicial
object.  These are standard consequences of Dold--Kan; see the Stacks
Project \cite{StacksNorm,StacksExact,StacksDK}.

Since $s_i$ inserts a loop in one of positions $2,\dots,n$, normalization
of the quotient $Q_n:=V_n/L_n$ gives
\begin{equation}\label{eq:normalizedquotient}
N(\Qcal)_{n-1}
\cong
V_n\Big/
\left(L_n+\Span\{e_S:S\text{ has a loop in }\{2,\dots,n\}\}\right).
\end{equation}
A tile whose only loop is at position $1$ survives this quotient.

The support bound in \Cref{cor:supportbound} now makes the normalized
problem finite.  The surviving columns at label number $n$ are exactly the
loopless $n$-label tiles together with the $(n-1)$-label loopless tiles
inserted with their unique loop at position $1$.  The only projected parent
rows that can be nonzero come from loopless parents or parents with their
unique loop at position $1$.

\begin{proposition}[Exact normalized ranks]\label[proposition]{prop:normalizedranks}
The normalized quotients have the following exact finite presentations:
\begin{center}
\begin{tabular}{rrrrr}
\toprule
$n$ & surviving tile columns & projected parent rows & row rank &
$\dim N(\Qcal)_{n-1}$\\
\midrule
6 & 1  & 0   & 0  & 1\\
7 & 15 & 7   & 6  & 9\\
8 & 50 & 47  & 31 & 19\\
9 & 66 & 121 & 54 & 12\\
10& 35 & 141 & 35 & 0\\
11& 5  & 46  & 5  & 0\\
\bottomrule
\end{tabular}
\end{center}
The displayed maximal nonzero minors can be chosen with determinants
\[
1,-1,1,-1,-1,1
\]
in order, with the first entry denoting the determinant $1$ of the
$0\times0$ minor of the rank-zero matrix.  The six rows of the table
are certificates \cert{NORM-6}, \cert{NORM-7}, \cert{NORM-8},
\cert{NORM-9}, \cert{NORM-10}, and \cert{NORM-11}.
For $n\ge12$ there are no normalized tile generators
(\cert{SUPPORT-CUTOFF}).
Consequently
\begin{equation}\label{eq:NQvanish}
N(\Qcal)_p=0\qquad(p\ge9).
\end{equation}
\end{proposition}

\begin{proof}
The column and row sets are determined exhaustively from the eleven
membership cores and the $29$ parent cores.  Rows from parents carrying a
loop in positions $2,\dots,n$ vanish after projection because every facet
retains that loop.  The remaining rows are literal transports of the
finite parent prototypes.

All ranks are computed over $\mathbb Q$.  For $n=10$ the selected
$35\times35$ submatrix has determinant $-1$, and for $n=11$
the selected $5\times5$ submatrix has determinant $1$.
These give full column rank.  For every displayed rank, the supplement
also supplies exact rational reduced rows and pivot data; checking that
every reconstructed row lies in their span gives the upper bound, while
the selected minor gives the lower bound.  This includes the non-full-rank
cases $n=7,8,9$.  The independent construction and its shared arithmetic
dependencies are specified in \Cref{app:repro}.
Finally, every tile has at most ten nonloop labels by
\Cref{cor:supportbound}; for $n\ge12$, every tile therefore has a loop in
positions $2,\dots,n$ and dies in the normalized quotient.
\end{proof}

\begin{remark}
Only the two high-degree zero statements at $n=10,11$ and the support
cutoff are needed for the all-$n$ vanishing argument.  The lower normalized
dimensions are useful for the rank formula below but are not an
extrapolation hypothesis.
\end{remark}

\section{Uniform global generation}
\label{sec:global}

The finite base case at $n=9$ is the last ingredient.  It is obtained by a
complete exact indicator-kernel certificate using the $942$ all-$Z$ tile
sources at $n=9$ (\cert{N9-UNIVERSE}).  The $693\times942$ local relation
matrix has rank $470$ (\cert{N9-LOCAL-RANK}).  An exact
$472\times942$ open-neighborhood indicator matrix has rank $472$, with a
$472\times472$ minor of determinant $4$, and annihilates every local
relation (\cert{N9-INDICATOR-RANK}, \cert{N9-ANNIHILATION}).
A separate implementation reconstructs all $444{,}624$ membership
decisions from the $472$ exact target records, including the spurion
inverse (\cert{N9-INDICATORS}).  Therefore \cert{N9-KERNEL} gives
\begin{equation}\label{eq:basecase}
K_9=L_9.
\end{equation}
We isolate the resulting infinite statement.

\begin{theorem}[Uniform indicator generation]\label{thm:KL}
For every $n\ge6$,
\[
K_n=L_n.
\]
\end{theorem}

\begin{proof}
Let
\[
\Hcal_p=K_{p+1}/L_{p+1},\qquad
\Qcal_p=V_{p+1}/L_{p+1}.
\]
By \Cref{thm:simplicial}, $\Hcal\hookrightarrow\Qcal$ is a monomorphism of
simplicial real vector spaces.

First, $\Hcal_p=0$ for $p\le8$.  Indeed, \eqref{eq:basecase} gives the
claim at label number $9$.  For every smaller label number, insert zero
columns up to nine labels and then apply the corresponding face maps.  The
face--degeneracy identities give left inverses to the insertion, while
\Cref{thm:amalgamation} guarantees that the faces preserve $L$.  Hence any
nonzero class in $K_r/L_r$ for $r<9$ would produce a nonzero class in
$K_9/L_9$, a contradiction.

Apply the exact normalization functor.  Since
$\Hcal\hookrightarrow\Qcal$, exactness gives
\[
N(\Hcal)\hookrightarrow N(\Qcal).
\]
For $p\le8$ the source $\Hcal_p$ is already zero, so
$N(\Hcal)_p=0$.  For $p\ge9$, \eqref{eq:NQvanish} gives
$N(\Qcal)_p=0$, and hence $N(\Hcal)_p=0$ as well.  Thus
\[
N(\Hcal)=0
\]
in every degree.  Faithfulness of normalization (equivalently, Dold--Kan)
forces $\Hcal=0$.  Therefore $K_n/L_n=0$ for every $n\ge6$.
\end{proof}

\subsection{An all-\texorpdfstring{$n$}{n} quotient-rank formula}
The normalized decomposition of a simplicial vector space is indexed by
order-preserving (weakly increasing) surjections $[n-1]\twoheadrightarrow[p]$,
where $[p]=\{0,\dots,p\}$.  There are $\binom{n-1}{p}$ such surjections.
Using \Cref{prop:normalizedranks}, one obtains
\begin{corollary}\label[corollary]{cor:rankformula}
For every $n$,
\begin{equation}\label{eq:Qrank}
\dim(V_n/L_n)
=\binom{n-1}{5}+9\binom{n-1}{6}
+19\binom{n-1}{7}+12\binom{n-1}{8},
\end{equation}
where binomial coefficients outside their natural range are zero.
Consequently
\[
\rank L_n=\dim K_n
=|\U_n|-\dim(V_n/L_n),
\]
with $|\U_n|$ given by \eqref{eq:Uncount}.
\end{corollary}

\section{The Magic Number corollary}
\label{sec:magic}

\begin{proof}[Proof of \Cref{thm:main}]
Let $\mathcal T$ and $\mathcal T'$ be complete multiplicity-one all-$Z$
tilings.  By \eqref{eq:tilingdifference},
\[
\chi_{\mathcal T}-\chi_{\mathcal T'}\in K_n.
\]
By \Cref{thm:KL}, this difference lies in $L_n$.  The augmentation
functional
\[
\epsilon:V_n\to\R,\qquad
\epsilon\!\left(\sum_S c_Se_S\right)=\sum_S c_S
\]
vanishes on $L_n$ by \eqref{eq:augmentationzero}.  Therefore
\[
|\mathcal T|-|\mathcal T'|=0.
\]
All complete all-$Z$ tilings thus have the same number of tiles.

A standard BCFW tiling exists for every admissible $n$ and has the Narayana
cardinality
\[
\frac{1}{n-3}\binom{n-3}{2}\binom{n-3}{3}
\]
for $k=2$, by the $m=4$ BCFW tiling theorem and the established BCFW count
\cite{KWZ,BCFWCluster}.  Comparing an arbitrary tiling to a standard BCFW
tiling gives the stated formula.  The algebraic simplification to
$(n-3)(n-4)^2(n-5)/12$ is immediate.
\end{proof}

\begin{remark}[Scope]
The theorem concerns constant-coefficient all-$Z$ indicator relations and
complete all-$Z$ positroid tilings.  It does not assert that the full
fixed-$Z$ indicator kernel equals $L_n$ for every individual positive
$Z$, that the flip graph is connected, or that $L_n$ is spanned by actual
global tiling differences for arbitrary $n$.  It also makes no claim for
$k>2$.
\end{remark}

\section{Discussion}

The proof separates into a rank-two local classification and a global
simplicial argument.  The local step is strongly rank-specific: it uses
the description of $\Gr_{\ge0}(2,n)$ by cyclic parallel classes and a
finite classification of nine-dimensional parents.  By contrast, the
final mechanism---adjacent amalgamation, loop insertion, normalization,
and support-bounded vanishing---is formulated at the level of source-cell
relations and may admit higher-rank analogues.  Determining which part of
this structure survives for $k\ge3$ is a natural next problem toward the
full $m=4$ Magic Number Conjecture.

The proof also contrasts with the $m=2$ argument of
\cite{M2Magic}.  There, a Parke--Taylor weight gives a direct unit valuation
on tiles.  Here the augmentation functional becomes well-defined only
after the global relation theorem $K_n=L_n$.  The finite $n=8,9$ indicator
certificates do admit signed rational unit functionals
(\cert{N8-UNIT}, \cert{N9-UNIT}), but their specific
weights do not appear to depend only on simple tile profile data
(the finite probes are \cert{WEIGHT-PROBES}).  A
canonical $m=4$ weight formula, if one exists, would provide an alternative
route to higher rank.

\appendix

\section{Exact finite certificates}\label{app:certificates}

The supplement is a directory named \path{publication/}.  Paths in this
appendix are relative to that directory.  A certificate ID denotes a
finite mathematical claim, its literal data, and the checks needed to
validate it.  These are distinct from software regression tests.  The
complete claim-to-file mapping is \path{CERTIFICATE_INDEX.md}; the source
of each check is \path{verifier/checks.py}, with exact linear algebra in
\path{verifier/exact.py}.  From the supplement directory, one may run
\begin{verbatim}
python verify_all.py --certificate MEM-CORES
python verify_all.py
\end{verbatim}
The first command replays the indicated family and its required inputs;
replacing \texttt{MEM-CORES} by any ID below selects that family.  The
second runs the full release verification.  A failed check terminates
with a nonzero exit status.  Only a successful full replay may print
\texttt{ALL PUBLICATION-CRITICAL CERTIFICATES PASSED}.

Integers and rational numbers are stored exactly; a rational string
\texttt{a/b} means $a/b\in\mathbb Q$.  Matrix rows and columns have the
order given in their files.  Indices into arrays are zero-based, while
source labels and the suffixes of public IDs are one-based.  Determinant
signs always refer to the published row and column order.  Stable IDs,
rather than a display order chosen by a reader, identify the certificates.

\subsection{Membership catalogue and exclusions}
\label{cert:MEM-CORES}\label{cert:MEM-FOUR-MASS}\label{cert:MEM-SPURION}

\paragraph{\texttt{MEM-CORES}.}
The literal catalogue is \path{certificates/membership/catalogue.json}.
Each record gives an ID \texttt{M001}, etc., the nonloop support $q$, the
number $s$ of ordered cyclic parallel classes, the classes themselves,
their sizes, the nonzero bases, the affine-permutation label, dihedral
orbit representative and stabilizer, membership status, and the applicable
mechanism.  Excluded cores are retained.  The verifier independently
enumerates the loopless rank-two cells with $q+s=12$, compares the full
labeled sets and dihedral data, and applies the criterion in
\Cref{thm:membership}.  The counts of admitted labeled cells are exactly
$1,14,36,30,5$ for $q=6,\dots,10$.  Choosing the nonloop support in
$[n]$ then gives \eqref{eq:Uncount}.  This check verifies the finite
exhaustion and arithmetic; the proof of the classification mechanisms is
in \Cref{thm:membership} and its cited input.

\paragraph{\texttt{MEM-FOUR-MASS}.}
\path{certificates/membership/four_mass.json} contains the exact matrices
$Z,C_1,C_2,G$ and the four-mass class pattern.  The check verifies every
increasing six-minor of $Z$ is positive; every source two-minor is
nonnegative with exactly the declared zero pattern; the stacked matrix
with row blocks $C_1,C_2$ has rank greater than two; and
\[
\det G\ne0,\qquad C_1Z=G(C_2Z).
\]
Thus $C_1,C_2$ represent distinct source planes in the same open positroid
cell and have the same image plane for this positive $Z$.  One such
counterexample suffices to exclude all-$Z$ injectivity.  The original
certificate is retained in
\path{certificates/membership/four_mass_historical.json}.

\paragraph{\texttt{MEM-SPURION}.}
\path{certificates/membership/spurion_inverse.json} contains exact source,
external, and target matrices for the inverse fixtures, together with
expected outcomes and their declared scope.  The verifier checks external
positivity, the source-image identities, and the exact inverse outcomes.
The generic inverse used in the local and $n=9$ certificates is specified
next.  Its finite fixtures are not an all-$Z$ injectivity proof; the
published spurion assertion cited in \Cref{thm:membership} is the uniform
geometric input.

\subsection{Exact membership decisions and open witnesses}
\label{app:inverse}

The implementation is \path{independent/membership.py}.  Its input is an
exact rational positive $n\times6$ matrix $Z$, a rank-two target frame
$Y$, and an admissible source label $S$.  The function \texttt{prepare}
removes loops and chooses a signed cyclic chart.  The decision function
\texttt{decide} returns $1$ exactly for $Y\in f_Z(S)$ under the generic
conditions below.  It returns $0$ for a strictly excluded target; an
exceptional equality or rank condition causes an error, not a silent
nonmembership answer.  These are decisions for the open source image.
The generic conditions ensure the same indicator value holds on an open
target neighborhood, as required for relations modulo dimension seven.

\paragraph{Complement-five inverse.}
The distinguished class $D$ is placed last.  The five complementary rows
span $H$, and the routine \texttt{complement} computes the line $Y\cap H$
in the basis of the first six external rows.  It requires $Y\not\subset H$
and all five line coordinates to be nonzero.  Mixed signs exclude a
positive source.  Otherwise their common sign is removed and a common
scale fixed, giving $a_i>0$.  For each complementary class $A_j$, form
\[
V_j=\sum_{i\in A_j}a_iZ_i.
\]
Choose a second target row with positive $h$-value and expand it in the
six-vector basis
\[
(V_1,\dots,V_\ell,Z_{D_1},\dots,Z_{D_r})
\]
from the proof of
\Cref{thm:membership}.  The coefficients give slopes $t_j$ and entries
$d_h$.  The exact source reconstructed from these coefficients is checked
against the target up to its explicitly computed invertible row change.
The routine requires the consecutive differences $t_{j+1}-t_j$ and all
$d_h$ to be nonzero.  It tests the signs of all source two-minors and,
when admitted, their precise zero/nonzero pattern.  Equivalently in this
chart, admission means $t_1<\cdots<t_\ell$ and every $d_h>0$.
A negative minor of the unique candidate excludes membership.  The inverse
argument in \Cref{thm:membership} proves necessity and sufficiency of
these checks, rather than merely giving a sufficient numerical test.

\paragraph{Spurion inverse.}
In the signed chart with three consecutive classes $A,B,D$ of size three,
write a candidate as
\[
C(x)=
\begin{pmatrix}
x_1&x_2&x_3&x_4&x_5&x_6&0&0&0\\
0&0&0&x_4&x_5&x_6&x_7&x_8&x_9
\end{pmatrix}.
\]
The routine \texttt{spurion} solves the twelve homogeneous ambient
coordinates of $C(x)Z=GY$ in the nine $x_i$ and four entries of $G$.
On the declared generic set, this $12\times13$ system has a
one-dimensional nullspace, each $x_i$ is nonzero, and the stack of $Y$
with each external class triple has rank five.  The check admits the
target exactly when the nine $x_i$ have a common sign, which may be
reversed.  It also verifies equality of image planes.  Indeed, three
distinct cyclic source directions can be put in the displayed gauge,
and positivity is precisely the common-sign condition.  The unique
nullspace therefore supplies both an inverse test and a nonmembership
test on this generic set.  The ambient formulation shares no production
matrix-building routine with the earlier quotient-coordinate inverse.

\paragraph{Openness.}
The recorded witnesses supply rational target frames and positive source
preimages where membership is asserted.  The replay checks the relevant
nonzero minors, ranks, and strict coefficient signs just described.
For the hostile and $n=9$ witnesses it also checks all increasing
four-twistors $\det(Y_1,Y_2,Z_i,Z_j,Z_k,Z_l)$ are nonzero.
Nonzero determinants and strict inequalities persist on an open
neighborhood.  Each matrix entry therefore records a locally constant
indicator value; a boundary sample is rejected.  This is why these rows
constrain $K_n$, whose definition ignores lower-dimensional sets.

\subsection{Parent cores and local relations}
\label{cert:LOCAL-PARENTS}\label{cert:LOCAL-RELATIONS}
\label{cert:LOCAL-TEN-WALL}\label{cert:LOCAL-HOSTILE-8}
\label{cert:LOCAL-HOSTILE-9}

\paragraph{\texttt{LOCAL-PARENTS} and \texttt{LOCAL-RELATIONS}.}
\path{certificates/local/parents.json} and its readable companion
\path{certificates/local/PARENTS.md} give the $29$ cores
\texttt{P001},\dots,\texttt{P029}.  Each records its ordered classes,
facet IDs, facet membership mechanisms, sector, relation dimension,
and full primitive coefficient vector in the published facet order.
Positive and negative facet-ID lists give the literal sides of every
$3\leftrightarrow3$ or $5\leftrightarrow5$ relation.
The verifier enumerates the types with $q+s=13$, reconstructs their
facets, compares every facet label, checks the partition $12+10+5+2=29$,
and verifies every coefficient sum is zero.  The original signed vectors
are retained in \path{certificates/local/relations_historical.json}.
The geometric proof in \Cref{thm:local} establishes that these vectors
are all-$Z$ relations; finite sign-vector arithmetic alone does not do so.

\paragraph{\texttt{LOCAL-TEN-WALL}.}
\path{certificates/local/overlap_witnesses.json} gives the $45$ witnesses,
with parent and witness IDs.  The five original family records, including
ordered walls and edge matrices, are
\path{certificates/local/overlaps_historical.json}.  Each witness gives
$Z,Y$, two exact positive preimages, source coefficients, and image row
changes.  The verifier checks these identities and uses the exact inverse
on all ten facets, including the eight asserted nonmembers.  Its row
has ones at the two present facets and zeros elsewhere, yielding the
coefficient equation $c_i+c_j=0$.  The nine rows form a $9\times10$
matrix; deleting the specified vertex column gives a minor of absolute
determinant one.  The known balanced sign vector lies in its kernel.
Consequently the kernel is one-dimensional.  This is the finite
upper-bound certificate for each ten-wall parent relation space.

\paragraph{\texttt{LOCAL-HOSTILE-8} and \texttt{LOCAL-HOSTILE-9}.}
The literal files \path{certificates/local/hostile_q8.json} and
\path{certificates/local/hostile_q9.json} contain the ordered eight
admissible facets, external matrix, eight target records, and indicator
matrix for each hostile parent.  Replaying all $64$ decisions in each
file reconstructs its $8\times8$ matrix; the exact determinant in the
published order must equal $1$.  This excludes every nonzero coefficient
relation on those eight facets.

\subsection{The complete \texorpdfstring{$n=9$}{n=9} certificate}
\label{cert:N9-UNIVERSE}\label{cert:N9-LOCAL-RANK}
\label{cert:N9-INDICATORS}\label{cert:N9-INDICATOR-RANK}
\label{cert:N9-ANNIHILATION}\label{cert:N9-KERNEL}

\paragraph{Universe and local rank.}
\texttt{N9-UNIVERSE} is
\path{certificates/n9/tiles.json}: the complete ordered list
\texttt{N9-T0001},\dots,\texttt{N9-T0942}.  The independent enumeration
must equal this list exactly, not merely have the same size.
\texttt{N9-LOCAL-RANK} is \path{certificates/n9/R9.json}, whose
$693\times942$ integer matrix $R_9$ is reconstructed from the transported
local prototypes.  The file includes row and column IDs, selected
independent rows, pivot columns, and sparse rational reduced rows.
The selected $470\times470$ minor has determinant $-1$, proving rank at
least $470$.  For the upper bound, write $B$ for the $470$ published
reduced rows and $p_1,\dots,p_{470}$ for their pivot positions.
The verifier checks $B_{i,p_j}=\delta_{ij}$ and, for every original row
$r$, the exact identity
\[
r=\sum_{j=1}^{470}r_{p_j}B_j.
\]
Thus every row lies in a $470$-dimensional space and the rank is exactly
$470$.  A nonzero minor alone would not prove this upper bound.

\paragraph{Witnesses and indicator rank.}
\texttt{N9-INDICATORS} uses
\path{certificates/n9/witnesses.json}, with a common exact external matrix
and target IDs \texttt{N9-W0001},\dots,\texttt{N9-W0472}.
The source frames, target frames, source labels, and recorded indicator
rows suffice to reconstruct every entry by \Cref{app:inverse}.
The independently rebuilt $472\times942$ binary matrix must agree
entrywise with both \path{certificates/n9/A9.json} and the original
\path{certificates/n9/indicator_historical.json}.  All
$472\cdot942=444{,}624$ decisions are replayed; none is inferred from
agreement on a sample.

For \texttt{N9-INDICATOR-RANK}, \path{certificates/n9/A9.json}
identifies the $472$ selected column indices in their determinant order;
its row IDs identify all $472$ selected rows.  The verifier extracts that
literal submatrix from the reconstructed matrix and checks its determinant
is $4$.  The lower bound is therefore $472$, and the number of rows is
the matching upper bound.

\paragraph{Annihilation and conclusion.}
\texttt{N9-ANNIHILATION} checks every entry of $A_9R_9^{\mathsf T}$ is
zero by exact integer arithmetic.  \texttt{N9-KERNEL} combines the
preceding checks with the geometric inclusion $L_9\subseteq K_9$ and
the open-witness implication $K_9\subseteq\ker A_9$:
\[
\dim L_9=470,\qquad \dim\ker A_9=942-472=470,
\qquad L_9\subseteq K_9\subseteq\ker A_9.
\]
It follows that $K_9=L_9$.  The finite matrix kernel is used as an upper
bound on the all-$Z$ relation space; no extrapolation from one $Z$ to all
positive $Z$ is made.

\subsection{Normalized quotients and support cutoff}
\label{cert:NORM-6}\label{cert:NORM-7}\label{cert:NORM-8}
\label{cert:NORM-9}\label{cert:NORM-10}\label{cert:NORM-11}
\label{cert:SUPPORT-CUTOFF}

The six literal files
\path{certificates/normalized/N6.json},\dots,
\path{certificates/normalized/N11.json} are \texttt{NORM-6},\dots,
\texttt{NORM-11}.  Each contains canonical column IDs and source labels,
row IDs and parent labels, the projected integer matrix, selected rows and
pivot columns, sparse rational reduced rows, and the exact selected
minor determinant.  The construction retains exactly the tile columns
with no loop in positions $2,\dots,n$.  Each transported parent relation
is projected onto those columns, with repeated facet images combined.
The reconstructed rows and labels must agree with the file in full.

The six dimensions, row counts, ranks, and quotient dimensions must equal
\Cref{prop:normalizedranks}.  The verifier uses the same exact row-span
upper certificate and selected-minor lower certificate as for $R_9$.
In particular, upper certificates are required at $n=7,8,9$.
At $n=10$ the $35\times35$ selected minor has determinant $-1$;
at $n=11$ the $5\times5$ selected minor has determinant $1$.
At $n=6$ there are no rows and the rank is zero; the empty determinant is
$1$ by convention and asserts no positive rank.

\texttt{SUPPORT-CUTOFF} ties the catalogue support $q\le10$ to the
normalization convention and checks that the normalized column set is
empty at $n=12$.  The statement for all $n\ge12$ is the symbolic
pigeonhole argument of \Cref{prop:normalizedranks}: with at most ten
nonloops there must be a loop in positions $2,\dots,n$.
The label number $n$ always corresponds to simplicial degree $p=n-1$.

\subsection{Additional finite observations}
\label{cert:N8-UNIT}\label{cert:N9-UNIT}\label{cert:WEIGHT-PROBES}

The Discussion's finite unit-function observations have separate IDs
\path{N8-UNIT}, \path{N9-UNIT}, and \path{WEIGHT-PROBES} in
\path{CERTIFICATE_INDEX.md}.  The first two expose the original rational
weights $\lambda$ and indicator matrices $A$ and check
$A^{\mathsf T}\lambda=\mathbf1$ exactly.  The third exposes the recorded
finite comparisons of weights with tile-profile data.  These are
supplementary finite observations, not premises of the all-$n$ theorem
or a claim of nonexistence of a general weight formula.

\section{Reproducibility and dependency map}\label{app:repro}

\Cref{tab:dependencies} separates the finite checks from the uniform
geometric arguments.  The supplement's \path{PROOF_DEPENDENCIES.md}
records the finer claim-to-certificate links.  An edge labelled mixed
uses both a symbolic argument and a finite certificate; numerical replay
does not verify the symbolic argument or an external theorem.

\begin{table}[ht]
\centering\small
\begin{tabular}{p{.27\textwidth}p{.33\textwidth}p{.31\textwidth}}
\toprule
Input & Consequence & Dependency type\\
\midrule
Membership mechanisms and \texttt{MEM-CORES}
 & Support bound and parent-facet classification & Mixed; spurion input external\\
\texttt{LOCAL-PARENTS}, \texttt{LOCAL-RELATIONS}, overlap and hostile certificates
 & Local theorem and $L_n\subseteq K_n$ & Mixed\\
Amalgamation, zero insertion, and truncation
 & Simplicial $V,L,K,K/L$ & Human/symbolic\\
Local theorem and the \texttt{N9-*} certificates
 & $K_9=L_9$ & Mixed\\
\texttt{NORM-10}, \texttt{NORM-11}, support bound
 & $N(\Qcal)_p=0$ for $p\ge9$ & Mixed\\
Simplicial structure, base case, normalized vanishing, Dold--Kan
 & $K_n=L_n$ for every $n\ge6$ & Human/symbolic and external citation\\
Balanced local vectors, global generation, standard BCFW tiling
 & Magic Number cardinality & Mixed and external citation\\
\bottomrule
\end{tabular}
\caption{Proof dependencies.  Monospaced certificate IDs mark the
machine-assisted inputs; the finite lower normalized ranks additionally
give \Cref{cor:rankformula}.}
\label{tab:dependencies}
\end{table}

\paragraph{What independent reconstruction means here.}
The historical generators and their outputs remain in the research
repository.  The release copies retain provenance rather than replacing
those artifacts.  The standalone reconstructor in \path{independent/} does not
import the production enumeration, facet, transport, or membership
functions from \path{src/}.  It reconstructs the finite labeled sets,
projected matrices, and witness decisions from the published conventions
and certificate inputs.  Its membership component adapts the historical
secondary ambient inverse verifier; it is not a third independently
invented inverse.  The release catalogue normalizer also shares the
standalone combinatorial routines, so independence refers to comparison
with the historical producers.  The implementations share the mathematical
classification and signed local prototypes, and use common exact
arithmetic libraries, including SymPy.  Their agreement is an independent
implementation check, not a separately proved classification theorem or
an independently implemented computer algebra system.  The supplement
distinguishes this reconstruction from replay of existing literal data
and from software regression tests.

\paragraph{Reproducibility files.}
\path{REPRODUCIBILITY.md} records the tested interpreter and platform,
pinned dependencies, installation and clean replay commands, resource
requirements, and cache behavior.  Verification uses exact arithmetic
and requires no network after dependency installation.  The manifest
\path{SHA256SUMS.txt} fixes the manuscript, PDF, literal certificates,
verifier sources, and release records.  \path{SUBMISSION_READINESS.md}
reports the actual release-validation state; the candidate designation
here does not assert that an archive has been deposited.
The historical symbolic notes are included under \path{proofs/}, notably
\path{proofs/research/run4/LOCAL_PARENT_THEOREM.md},
\path{proofs/research/run6/MEMBERSHIP_CORES.md},
\path{proofs/research/run6/SPURION_INVERSION.md},
\path{proofs/research/run6/N9_INDICATOR_KERNEL.md}, and the run7
amalgamation, simplicial, and global-generation notes.

\begingroup
\emergencystretch=2em
\paragraph{Code and data availability.}
The accompanying supplement is the candidate release
\texttt{v1.0.0-paper}.  Its historical base commit is
\path{284c19a88a8e1327f139d4e3c58d16c751610b73}; a release commit has
not yet been minted.  The intended public repository is
\textbf{[PUBLIC REPOSITORY URL TO BE INSERTED]}, and the archival
identifier is \textbf{[ZENODO DOI TO BE INSERTED AFTER RELEASE]}.
Neither placeholder denotes an existing upload.  From the deposited
supplement's root, the verification command is
\texttt{python verify\_all.py}.  Code is offered under the MIT license;
the manuscript and data are offered under CC BY 4.0.
\par\endgroup

\section*{AI assistance disclosure}
OpenAI ChatGPT and Codex models assisted with proof exploration,
finite-case formalization, code generation, verification tooling,
adversarial proof review, and manuscript drafting and editing.  AI systems
are not authors.  The human author remains responsible for all
mathematical claims, computations, citations, and the final manuscript.

\begin{thebibliography}{99}

\bibitem{Postnikov}
A.~Postnikov,
\emph{Total positivity, Grassmannians, and networks},
arXiv:math/0609764 (2006).

\bibitem{ABCGPT}
N.~Arkani-Hamed, J.~L.~Bourjaily, F.~Cachazo, A.~B.~Goncharov,
A.~Postnikov, and J.~Trnka,
\emph{Scattering amplitudes and the positive Grassmannian},
arXiv:1212.5605 (2012).

\bibitem{GalashinLam}
P.~Galashin and T.~Lam,
\emph{Parity duality for the amplituhedron},
Compositio Mathematica \textbf{156} (2020), 2207--2262,
\href{https://doi.org/10.1112/S0010437X20007411}{doi:10.1112/S0010437X20007411}.

\bibitem{KWZ}
\begingroup\raggedright
S.~N.~Karp, L.~K.~Williams, and Y.~X.~Zhang,
\emph{Decompositions of amplituhedra},
Annales de l'Institut Henri Poincar\'e D \textbf{7} (2020), 303--363,
\href{https://doi.org/10.4171/AIHPD/87}{doi:10.4171/AIHPD/87}.
\par\endgroup

\bibitem{BCFWCluster}
C.~Even-Zohar, T.~Lakrec, M.~Parisi, R.~Tessler,
M.~Sherman-Bennett, and L.~Williams,
\emph{Cluster algebras and tilings for the $m=4$ amplituhedron},
arXiv:2310.17727v3 (August 12, 2026), Theorems 8.7 and 13.3.

\bibitem{ClusterResults}
C.~Even-Zohar, T.~Lakrec, M.~Parisi, R.~Tessler,
M.~Sherman-Bennett, and L.~Williams,
\emph{A cluster of results on amplituhedron tiles},
Letters in Mathematical Physics \textbf{114}, article 111 (2024),
\href{https://doi.org/10.1007/s11005-024-01854-4}{doi:10.1007/s11005-024-01854-4}.

\bibitem{M2Magic}
M.~Parisi, M.~Sherman-Bennett, R.~Tessler, and L.~Williams,
\emph{The Magic Number Conjecture for the $m=2$ amplituhedron and
Parke--Taylor identities},
arXiv:2404.03026 (2024).

\bibitem{StacksNorm}
The Stacks Project,
\emph{Simplicial objects: quotient by degeneracies},
Tag 019A, \url{https://stacks.math.columbia.edu/tag/019A}.

\bibitem{StacksExact}
The Stacks Project,
\emph{Simplicial objects: exactness of normalization},
Tag 019B, \url{https://stacks.math.columbia.edu/tag/019B}.

\bibitem{StacksDK}
The Stacks Project,
\emph{Dold--Kan},
Tag 019D, \url{https://stacks.math.columbia.edu/tag/019D}.

\end{thebibliography}

\end{document}
